\begin{lemma}[Reflecting a running minimum]\label{fprof:reflection}
Let \(0<a\leq1/4\), \(a\leq b\leq1\), and let
\(g:[0,N]\to\mathbb R\) be 1-Lipschitz with
\(ax\leq g(x)\leq bx\). For \(3N/4\leq n_0<N\), there is an
admissible decreasing chain from \(n_0\) through \(N/2\) to zero,
with \(O(1+\log(N/(N-n_0)))\) edges and total cost at most
$$
 N\max\{E(a,b),M(a,b)\},\qquad
 E(a,b)=\frac{1-2a-2a^2+(2+a)b}{6(1+a)},\qquad
 M(a,b)=\frac{3b}{8}+\frac3{16}-\frac{3a}{4}.
$$
On a grid of mesh \(T\) containing \(N/2,3N/4,n_0,N\), the endpoints
can be chosen on the grid with an additional \(O(KT)\) cost, where \(K\)
is the displayed length bound. If the barriers hold with additive error
\(e_0N\), the additional cost is \(O(Ke_0N+KT)\).
\end{lemma}
\begin{proof}
Scale to \(N=1\). Write \(v=g(1/2)\), \(e=g(1)\), and choose
a minimum \(t\) of \(g\) on \([1/2,1]\), with \(m=g(t)\).
If \(t\leq3/4\), the early-minimum construction in
Lemma~\ref{fprof:split} gives a chain of cost at most \(E(a,b)\).
Its proof uses the cone above the minimum to give
$$
 v-m+\frac{1-t+m-e}{3};
$$
maximizing this under
\(e\geq a\), \(m\geq at\), and
\(v\leq\min\{b/2,m+t-1/2\}\) gives precisely \(E(a,b)\).

It remains to treat \(t>3/4\). Choose \(u\in[1/2,3/4]\) at which
\(g\) attains its minimum on that interval, and put \(w=g(u)\).
In particular,
\begin{equation}\label{fprof:headconstraint}
 w\geq v-u+\tfrac12,\qquad w\geq m.
\end{equation}
We give two chains, retaining the information about \(u,w\) in both bounds.

Define the nondecreasing running deficit
$$
 r(x)=\begin{cases}
 0,&x\leq u,\\
 w-\min_{z\in[u,x]}g(z),&x\geq u,
 \end{cases}
 \qquad h(x)=g(x)+r(x).
$$
The second expression for \(r\) is nonnegative because \(g(u)=w\).
The function \(h\) is 1-Lipschitz. Indeed, on \([u,1]\) it is at least
\(w\). For \(u\leq x<y\), if \(r(x)=r(y)\), its increment equals the
increment of \(g\). If \(r(y)>r(x)\), a new running minimum is attained
at some \(z\in[x,y]\), so
$$
 h(y)=g(y)+w-g(z),\qquad h(x)\geq w,
 \qquad h(y)-h(x)\leq y-x.
$$
The reverse inequality follows from
\(h(y)-h(x)\geq g(y)-g(x)\geq-(y-x)\).
Gluing to \(g\) at \(u\) proves the assertion on the whole interval.
Moreover,
$$
 \min_{[1/2,1]}h=h(u)=w,\qquad h(1/2)=v,\qquad
 h(1)=e+w-m.
$$
The cone construction for an early minimum applies to \(h\), and gives
a chain with cost at most
$$
 v-w+\frac{1-u+w-(e+w-m)}3.
$$
No upper linear barrier for \(h\) is required here: that construction
uses the upper cone \(h(z)\leq w+z-u\) for \(z\geq u\), its early
global minimum, and its endpoint value.

For every edge \(p<q\), monotonicity of \(r\) gives
$$
 \min_{[p,q]}h\leq\min_{[p,q]}g+r(q),\qquad
 c_g(p,q)\leq c_h(p,q)+r(q)-r(p).
$$
The deficit increments telescope along the chain, with total at most
\(r(1)=w-m\). Thus the same chain for \(g\) has cost at most
\begin{equation}\label{fprof:reflectionR}
 R_0=v+\frac{1-u-e-2m}{3}.
\end{equation}

For the second chain, first suppose \(t\leq n_0\). Use the upper cone
from the global minimum at \(t\) to reach \(t\), at cost at most
\((1-t+m-e)/3\). This is the same cone-potential construction, now stopped
at \(t\); it does not require \(t\leq3/4\). From \(t\), take maximal
admissible jumps toward \(3/4\), clipping the last jump at \(3/4\).
Each edge cost is at most its length, so their total cost is at most
\(t-3/4\). Finally the admissible edge from \(3/4\) to \(1/2\)
has cost exactly \(v-w\).

If \(t>n_0\), omit the first part and start the clipped jumps at \(n_0\).
Their cost is at most \(n_0-3/4\leq t-3/4\). The omitted cone bound is
nonnegative since \(e-m\leq1-t\). In either case this second chain has
cost at most
\begin{equation}\label{fprof:reflectionS}
 S_0=v-w+\frac{2t+m-e-5/4}{3}.
\end{equation}
The jump from \(1/2\) to zero has zero cost.

Both chains have \(O(1+\log(1/(1-n_0)))\) edges. This was established for
the cone construction in Lemma~\ref{fprof:split}; in the additional clipped
jumps, complementary depth doubles except on the last edge. If \(t\leq n_0\),
starting that part at \(t\) cannot increase the stated length bound.

Choose the cheaper of the two chains. By \eqref{fprof:headconstraint},
$$
 S_0\leq u+\frac{2t+m-e-11/4}{3}.
$$
Consequently its cost is at most
\begin{align*}
 \min\{R_0,S_0\}
 &\leq\frac34 R_0+\frac14 S_0\\
 &\leq\frac{3v}{4}+\frac1{48}-\frac e3+\frac t6-\frac{5m}{12}\\
 &\leq\frac{3b}{8}+\frac1{48}-\frac a3+\frac{2-5a}{12}\,t\\
 &\leq\frac{3b}{8}+\frac3{16}-\frac{3a}{4}=M(a,b).
\end{align*}
Here we used \(v\leq b/2\), \(e\geq a\), \(m\geq at\), \(t\leq1\),
and \(2-5a\geq0\). Combining the early and late cases proves the real bound.

For the grid assertion, round every endpoint of the continuous chain downward
to the nearest grid point, retaining the prescribed endpoints, which are already
on the grid. Write \(p',q'\) for the rounded versions of an admissible edge
\(p<q\). After dividing by the mesh, with \(N/T\) an integer,
$$
 \lfloor p/T\rfloor\geq\lfloor 2q/T-N/T\rfloor
 \geq 2\lfloor q/T\rfloor-N/T.
$$
Thus the rounded edge remains admissible. Delete duplicate consecutive endpoints.
An endpoint moves by at most \(T\), and the minimum of a 1-Lipschitz function
on an interval changes by at most \(T\) when both endpoints move by at most
\(T\). Therefore the cost changes by at most \(2T\) per edge. The forced
midpoint is preserved, and the edge count cannot increase. This proves the
\(O(KT)\) assertion without requiring a grid minimum for the reflected function.
Finally clamp an approximate-barrier profile
between \(ax\) and \(bx\), then interpolate grid values. Clamping preserves
the Lipschitz constant and changes the profile by \(O(e_0N+T)\); each edge
cost changes by at most twice this uniform error. This proves stability.
\end{proof}

\begin{remark}[The sufficient region]\label{fprof:reflectionregion}
For \(0<a\leq1/4\), put
$$
 b_E(a)=\frac{8a^2+8a-1}{a+2}.
$$
Then \(E(a,b)<a\) is equivalent to \(b<b_E(a)\), and
$$
 M(a,b_E(a))-a=\frac{5a(4a-1)}{16(a+2)}\leq0.
$$
Because \(M\) is strictly increasing in \(b\), the strict condition
\(b<b_E(a)\) also implies \(M(a,b)<a\).
Thus \(E(a,b)<a\) alone implies
\(\max\{E(a,b),M(a,b)\}<a\) in this parameter range.
This is a sufficient-condition statement; it does not assert a universal
chain bound by \(E(a,b)\) for all larger values of \(b\).
\end{remark}
